\section{Post-saturation successor: scientific authority composition over heterogeneous donor systems}

Modern agent systems already provide strong local authority mechanisms. The successor therefore treats evidence-backed permission graphs, proof-carrying actions, delegation provenance, deterministic authority narrowing, principal-chain authorization, typed verifier certificates, heterogeneous authorization-evidence chains, cross-domain authority relations, scientific release harnesses, claim-to-evidence verification, contract-governed research artifacts, evidence-calibrated claim adjudication, and evidence-ledger review as donor-owned authority objects. P8 asks a different compositional question: when may one of those local judgments discharge a target scientific obligation of a particular domain, kind, scope, content and epoch?

Let $a$ be a donor authority, verification, adjudication, or release object with native predicate $D(a)$. The target scientific-discharge type is
\[
\tau=(d,k,s,c,e),
\]
where the five coordinates denote domain, judgment kind, scope, content identity and epoch. The lifted state also records whether authority has only been preserved or narrowed along the relevant chain, a blocker state $b\in\{\mathrm{REFUTED},\mathrm{UNDETERMINED},\mathrm{ESTABLISHED}\}$, two alternative complete support families, and an optional protected type-coercion witness.

A donor-derived judgment may discharge the target only if $D(a)$ passes, the scientific type matches or is explicitly bridged by a complete protected coercion, scientific authority has not widened through the chain, the blocker is \textsc{refuted}, and at least one complete support family remains valid. An \textsc{undetermined} blocker yields \textsc{cannot\_check}; an \textsc{established} blocker blocks discharge.

\paragraph{Native-authority conservativity.}
The lifted calculus does not change the donor-native action, delegation, evidence, adjudication, or release verdict. A local certificate remains valid for its native object even when it does not authorize a broader scientific conclusion.

\paragraph{Typed scientific-discharge separation.}
For every registered donor family and every non-inert coordinate of $\tau$, a native-valid local judgment can fail to discharge a target that differs only on that coordinate. Local authority is therefore reusable but is not an unrestricted authority token.

\paragraph{Monotone heterogeneous propagation.}
Scientific authority may be preserved or narrowed through heterogeneous chains. A hop that widens the inherited scientific type cannot discharge the widened target merely because each local donor step is native-valid.

\paragraph{Protected coercion.}
A complete subject- and epoch-bound protected coercion may bridge an otherwise incompatible type. The matched judgment without that explicit bridge remains blocked. Intent similarity, semantic similarity, common action digest, claim fluency, or generic action composability is not by itself a coercion.

\paragraph{Blocker and revocation laws.}
Under otherwise satisfied conditions, \textsc{refuted} permits discharge, \textsc{undetermined} yields \textsc{cannot\_check}, and \textsc{established} blocks. When two independent complete support families establish the same target, revoking one family preserves discharge through the other; invalidating every complete support family removes it.

\paragraph{Heterogeneous receipt insufficiency boundary.}
Valid composition of heterogeneous authorization receipts, canonical action or subject binding, and successful local policy evaluation are useful prerequisites. They are not proof that a different scientific target is established. The target type, blocker state, support families and any scientific coercion remain explicit.

\paragraph{Permission, release and support are not inverse relations.}
Native denial of one action or release path does not scientifically refute a proposition that is independently established through another complete evidence family. Conversely, permission or release for one object does not discharge a broader or differently typed scientific obligation.

\paragraph{Ideal decentralized-product equivalence.}
A decentralized donor product equipped with exactly the same scientific type, narrowing, blocker, support-family, coercion and composition semantics agrees extensionally with P8. The contribution is the explicit scientific-discharge composition interface rather than centralized expressive power.

\paragraph{Exhaustive bounded support.}
The final X4 model spans thirteen donor authority families and 39,936 exact states. It contains zero donor-conservativity and ideal-product mismatches; 65 minimal scientific-type separation witnesses; 65 protected-coercion successes and 65 matched unprotected countermodels; the three-state blocker law for all thirteen donor families; 26 one-support-family revocation survivals and 13 all-support-family blocks; and 169 valid heterogeneous ordered-chain compositions with 169 matched scientific-authority-widening countermodels. An independent implementation reproduces the canonical row digest \texttt{ed186b824692fd5b3ab31be718c75b84e2126b577ce921ca5cc01b2d08ae19e6}.

The successor therefore widens P8 from a cross-capability anti-laundering interface into a constructive scientific-authority composition calculus over absorbed security, delegation and scientific-verification systems. The result remains bounded formal evidence; it does not establish deployed-agent superiority, a universal full-manuscript release gate, or inherent centralization advantage.
